%% =====================================================================
%%  T-10-04  The Attention Market
%%  -------------------------------------------------------------------
%%  One idea per file. Core is mandatory. Slots in canonical order.
%%  Max 700 words. No layout commands. No filler. New source material
%%  for this entry goes in sources/B3/T-10-04.tex -- never by
%%  rewriting this file (docs/RULES.md R0.1).
%%  =====================================================================
%% @kind: topic
%% @id: T-10-04
%% @title: The Attention Market
%% @subtitle: What is unseen counts for nothing --- attention must be claimed
%% @chapter: c10
%% @order: 40
%% @status: stable
%% @sources: B3
%% @updated: 2026-09-19

\topic{T-10-04}{The Attention Market}{What is unseen counts for nothing --- attention must be claimed}

\begin{core}
Quality without visibility is unfinished work: everything is judged by appearance, and what is unseen counts for nothing. B3's blunt instruction is to stand out, be conspicuous, make yourself a magnet of attention --- larger, more colourful, more mysterious than the crowd. This chapter handles absence and scarcity elsewhere; this entry is the affirmative side: attention is claimed, not merely merited.
\end{core}
\begin{mechanism}
Attention is the gateway resource: no door opens for the person the deciders cannot see. The source's premise is that judgement happens on appearances because appearance is all a busy audience can process; buried in oblivion, your substance is arithmetic without a market. Conspicuousness converts work into reputation (T-09-01) and reputation into inbound leverage (T-14-05). The mechanism has a two-sided shape --- this chapter's other entries sell absence, and the display itself must retain some mystery: the fully explained person is fully priced. The source's radical form, court attention at all costs, is a priority claim: when choosing between another increment of quality and another of visibility, most practitioners over-buy quality --- because quality is the part they control and the part their craft ethics reward them for.
\end{mechanism}
\begin{conditions}
Strongest in crowded fields where deciders skim --- hiring, publishing, promotion, markets with more sellers than attention. Weakens in small closed systems where everyone already knows everyone's work, and in contexts that punish standing out --- there the camouflage logic of T-18-04 governs instead. Fatal misuse: attention without substance. Visibility is an amplifier, and amplifiers amplify zeros.
\end{conditions}
\begin{application}
  \item Decide what you are knowable \emph{for}, first. Attention without an association is noise.
  \item Attach your name to the visible artefact --- the presentation, the byline, the room --- not the effort behind it.
  \item Court deciders, not crowds. Crowds admire; deciders allocate.
  \item Keep one element unexplained. Mystery keeps attention alive; total legibility ends it.
  \item Schedule visibility like work. The demo, the published note, the panel seat --- calendar it or it does not happen.
\end{application}
\begin{feedback}
  \item Working: strangers describe your work back to you accurately. It has entered the market's memory.
  \item Working: opportunities arrive addressed to the association, not posted into a pool.
  \item Failing: your best work is discovered post-hoc and credited to someone louder. The claim was never made.
  \item Failing: you are known for visibility itself. The association failed; the amplifier found no signal.
\end{feedback}
\begin{failure}
Attention is a debt as much as an asset: it invites scrutiny, envy (T-09-04) and rivals, and it must be paid for continuously --- an unfed audience drifts. Over-claiming without substance converts the association into a warning label. And the habit corrodes the craft: performers begin optimising for the notice rather than the thing noticed, a trade the field punishes on a delay everyone claims surprise at.
\end{failure}
\begin{countermeasures}
  \item Separate claim from substance when evaluating anyone: what do they make, what are they known for? The gap is the fact sheet.
  \item Watch the association trick: one true credit attached to a portfolio of borrowed ones. Audit the portfolio.
  \item In your own shop, occasionally reward the visible work of quiet builders --- the market's attention bias is real and you are one of its instruments.
  \item Notice who goes quiet before going loud. The cultivated absence before the entrance is this chapter's other half (T-10-05).
\end{countermeasures}

\sources{B3}
\seesources
\seealso{T-10-05, T-09-01, T-10-06}
